\section{Webapps}
\label{research:Webapps}
\textbf{Source:} \cite{asterWhatAreAdvantages2025} \cite{CapacitorCrossplatformNative}
At this point in development, I was certain that building a web app was the most practical choice for my project. Web apps offer significant advantages, such as ease of development and cross-platform compatibility, meaning the same application can be accessed from a wide range of devices, including desktops, tablets, and smartphones. With web technologies evolving rapidly, frameworks like React, Vue, and Angular make it easier than ever to create dynamic and responsive apps that run efficiently across different platforms. However, I was also aware that web apps have some inherent limitations, particularly when it comes to integrating with the native functionality of devices.

Despite these limitations, I wanted to explore ways I could bridge the gap and integrate more natively with devices. There are several tools and frameworks that facilitate this, such as Ionic Capacitor, which allows web apps to access native mobile features like the camera, GPS, and local storage. Capacitor makes it possible for web apps to act more like native applications, providing access to native APIs, such as for file handling and push notifications, without needing to write separate native code for each platform. Additionally, I considered integrating PWA (Progressive Web App) features, which allow web apps to be installed on a device and run offline, further blurring the line between web and native apps.

Another key area I wanted to explore was utilizing device-specific features like the file system, which is crucial for my app’s functionality. For example, I could enable users to upload and manage files directly from their devices, access local files, and perform actions like saving or retrieving documents. Tools like FileSystem API and Web Bluetooth API offer direct interaction with the device hardware, allowing me to integrate more advanced functionalities into the web app. Using these tools, I could access the device’s native camera to capture images or scan barcodes, and leverage sensors for geolocation or motion tracking, all without needing to develop separate applications for each platform.

However, while these capabilities are becoming more accessible on many platforms, there are challenges with iOS. Apple has historically restricted some of these native integrations in mobile Safari. For instance, the FileSystem API and certain sensor access features are either limited or unavailable in the iOS version of Safari, making it harder to fully integrate these features on iPhones and iPads. This has been a consistent issue for developers looking to create cross-platform apps that rely on these native functionalities. In contrast, Android devices tend to provide broader access to these APIs, allowing for a more seamless experience across devices.

Fortunately, for my current project, I don’t foresee these limitations significantly impacting the overall functionality of my app. While it’s true that integrating more native device capabilities on iOS will require creative workarounds, the core features I plan to implement should be accessible on most devices, including Apple products. However, this is an important consideration for the future, as the landscape of web app capabilities continues to evolve. As I move forward, I will keep a close eye on any changes in Apple’s policies and the capabilities of mobile web browsers to ensure my app remains fully functional and user-friendly across all platforms.
